% from ega1-9-fr.tex
\oldpage[I]{180}
Module de type fini~; or, il est isomorphe (en tant que
$\mathscr{B}$-Module) à $\widetilde{N}$, où $N$ est le noyau de
l'homomorphisme correspondant de $B$-modules $B^m\to B$
(\hyperref[I.1.3.13-fr]{1.3.13}). Comme $B$ est noethérien, le sous-module
$N$ de $B^m$ est un $B$-module de type fini, donc il existe un homomorphisme
$B^p\to B^m$ d'image $N$~; la suite $B^p\to B^m\to B$ étant exacte, il
en est de même de la suite
$\mathscr{B}^p\to\mathscr{B}^m\to\mathscr{B}$ correspondante
(\hyperref[I.1.3.5-fr]{1.3.5}) et comme $\mathscr{N}$ est l'image de
$\mathscr{B}^p\to\mathscr{B}^m$
(\hyperref[I.1.3.9-fr]{1.3.9, (i)}), la proposition est démontrée.
\end{proof}

\begin{corollary}[9.6.4]
\label{I.9.6.4-fr}
Sous les hypothèses de (\hyperref[I.9.6.3-fr]{9.6.3}), pour qu'un
$\mathscr{B}$-Module $\mathscr{F}$ soit cohérent, il faut et il suffit
qu'il soit un $\mathscr{O}_X$-Module quasi-cohérent et un
$\mathscr{B}$-Module de type fini. S'il en est ainsi, et si $\mathscr{G}$
est un sous-$\mathscr{B}$-Module ou un $\mathscr{B}$-Module quotient de
$\mathscr{F}$, pour que $\mathscr{G}$ soit un $\mathscr{B}$-Module
cohérent, il faut et il suffit que $\mathscr{G}$ soit un
$\mathscr{O}_X$-Module quasi-cohérent.
\end{corollary}

\begin{proof}
Compte tenu de (\hyperref[I.9.6.1-fr]{9.6.1}), les conditions sur
$\mathscr{F}$ sont évidemment nécessaires~; montrons qu'elles sont
suffisantes. On peut se limiter au cas où $X$ est affine d'anneau
noethérien $A$, $\mathscr{B}=\widetilde{B}$, où $B$ est une $A$-algèbre de
type fini, $\mathscr{F}=\widetilde{M}$, où $M$ est un $B$-module, et où il
existe un $\mathscr{B}$-homomorphisme surjectif
$\mathscr{B}^m\to\mathscr{F}\to 0$. On a alors une suite exacte
correspondante $B^m\to M\to 0$, donc $M$ est un $B$-module de type fini~;
en outre, le noyau $P$ de l'homomorphisme $B^m\to M$ est alors un
$B$-module de type fini, puisque $B$ est noethérien. On en conclut
(\hyperref[I.1.3.13-fr]{1.3.13}) que $\mathscr{F}$ est le conoyau d'un
$\mathscr{B}$-homomorphisme $\mathscr{B}^n\to\mathscr{B}^m$, et est donc
cohérent puisque $\mathscr{B}$ est un faisceau cohérent d'anneaux
(\hyperref[0.5.3.4-fr]{0, 5.3.4}). Le même raisonnement montre qu'un
sous-$\mathscr{B}$-Module (resp. un $\mathscr{B}$-Module quotient)
quasi-cohérent de $\mathscr{F}$ est de type fini, d'où la seconde partie du
corollaire.
\end{proof}

\begin{proposition}[9.6.5]
\label{I.9.6.5-fr}
Soit $X$ un schéma quasi-compact ou un préschéma dont l'espace sous-jacent
est noethérien. Pour toute $\mathscr{O}_X$-Algèbre quasi-cohérente de type
fini $\mathscr{B}$, il existe un sous-$\mathscr{O}_X$-Module quasi-cohérent
de type fini $\mathscr{E}$ de $\mathscr{B}$ tel que $\mathscr{E}$ engendre
(\hyperref[0.4.1.4-fr]{0, 4.1.4}) la $\mathscr{O}_X$-Algèbre
$\mathscr{B}$.
\end{proposition}

\begin{proof}
En effet, il existe par hypothèse un recouvrement fini $(U_i)$ de $X$ formé
d'ouverts affines tels que $\Gamma(U_i,\mathscr{B})=B_i$ soit une algèbre
de type fini sur $\Gamma(U_i,\mathscr{O}_X)=A_i$~; soit $E_i$ un
sous-$A_i$-module de type fini de $B_i$ engendrant la $A_i$-algèbre $B_i$~;
en vertu de (\hyperref[I.9.4.7-fr]{9.4.7}), il existe un
sous-$\mathscr{O}_X$-Module $\mathscr{E}_i$ de $\mathscr{B}$,
quasi-cohérent et de type fini, tel que
$\mathscr{E}_i|U_i=\widetilde{E_i}$. Il est clair que la somme
$\mathscr{E}$ des $\mathscr{E}_i$ répond à la question.
\end{proof}

\begin{proposition}[9.6.6]
\label{I.9.6.6-fr}
Soit $X$ un préschéma dont l'espace sous-jacent est localement noethérien,
ou un schéma quasi-compact. Alors toute $\mathscr{O}_X$-Algèbre
quasi-cohérente $\mathscr{B}$ est limite inductive de ses
sous-$\mathscr{O}_X$-Algèbres quasi-cohérentes de type fini.
\end{proposition}

\begin{proof}
En effet, il résulte de (\hyperref[I.9.4.9-fr]{9.4.9}) que $\mathscr{B}$
est limite inductive (en tant que $\mathscr{O}_X$-Module) de ses
sous-$\mathscr{O}_X$-Modules quasi-cohérents de type fini~; ces derniers
engendrent des sous-$\mathscr{O}_X$-Algèbres quasi-cohérentes de type fini
de $\mathscr{B}$ (\hyperref[I.1.3.14-fr]{1.3.14}), dont $\mathscr{B}$ est
\emph{a fortiori} limite inductive.
\end{proof}

% (continued in ega1-10-fr.tex, before its first page marker)
\section{Schémas formels}
\label{section:I.10-fr}

\subsection{Schémas formels affines.}
\label{subsection:I.10.1-fr}

\begin{env}[10.1.1]
\label{I.10.1.1-fr}
Soit $A$ un anneau topologique \emph{admissible}
(\hyperref[0.7.1.2-fr]{0, 7.1.2})~; pour tout idéal de définition
$\mathfrak{J}$ de $A$, $\operatorname{Spec}(A/\mathfrak{J})$ s'identifie
au sous-espace fermé $V(\mathfrak{J})$ de $\operatorname{Spec}(A)$
(\hyperref[I.1.1.11-fr]{1.1.11}), ensemble des idéaux premiers
\emph{ouverts} de $A$~; cet espace topologique ne dépend pas de
